\section{Полный контекст подтверждённых эффектов}
\subsection{Все восемь устойчивых gain-only взаимодействий CLINC150}
\begin{table}[htbp]\centering\small
\caption{Все взаимодействия gain-методов, обнаруженные при
accuracy-настройке в исходной семье C и сохранённые после 10\,000
повторов. Ориентация выбрана так, чтобы изменение разности было
положительным. Строки зависимы и относятся к трём парам экспертов.}
\label{tab:all_eight}
\begin{tabular}{llrrr}
\toprule
Пара & Методы & $\Delta$, п.п. & 95\% CI & $p_{adj}$ \\
\midrule
TF3 $\to$ TF10 & Hurdle--Soft-gain & 0{,}77 & [0{,}02; 1{,}53] & 0{,}0401 \\
TF10 $\to$ TFfull & Events--Direct & 0{,}69 & [0{,}05; 1{,}34] & 0{,}0264 \\
TF10 $\to$ TFfull & Ternary--Direct & 0{,}77 & [0{,}11; 1{,}43] & 0{,}0112 \\
TF10 $\to$ TFfull & Joint--Direct & 0{,}71 & [0{,}02; 1{,}39] & 0{,}0382 \\
TF10 $\to$ TFfull & Ternary--Soft-gain & 0{,}79 & [0{,}05; 1{,}53] & 0{,}0264 \\
TFfull $\to$ MiniLM & Hurdle--Direct & 0{,}87 & [0{,}06; 1{,}68] & 0{,}0264 \\
TFfull $\to$ MiniLM & Hurdle--Events & 0{,}68 & [0{,}03; 1{,}33] & 0{,}0334 \\
TFfull $\to$ MiniLM & Hurdle--Soft-gain & 0{,}81 & [0{,}07; 1{,}54] & 0{,}0211 \\
\bottomrule
\end{tabular}

\end{table}

\begin{table}[htbp]\centering\footnotesize
\caption{Численная устойчивость основных CLINC-контрастов.
Первые интервалы --- исходный расчёт 2\,000 повторов,
вторые --- однократное уточнение до 10\,000; p-values относятся
к уточнённому расчёту. Семья из 1\,698 контрастов одинакова.}
\label{tab:mc_audit}
\begin{tabular}{lrrr}
\toprule
Контраст & CI: 2\,000 & CI: 10\,000 & $p_{adj}$ \\
\midrule
Direct--Weak-error, interaction & [-3{,}34; -0{,}66] & [-3{,}34; -0{,}66] & 0{,}0005 \\
Hurdle--Direct, harm8 & [0{,}04; 1{,}65] & [0{,}03; 1{,}66] & 0{,}0351 \\
Hurdle--Direct, interaction & [0{,}06; 1{,}67] & [0{,}06; 1{,}68] & 0{,}0264 \\
Ternary--Direct, interaction & [0{,}10; 1{,}44] & [0{,}11; 1{,}43] & 0{,}0112 \\
Direct, IPW--raw & [0{,}68; 3{,}40] & [0{,}68; 3{,}40] & 0{,}0005 \\
\bottomrule
\end{tabular}

\end{table}

В полной семье C при уточнении изменились шесть граничных статусов:
один gain-only MSE-контраст стал значимым, три взаимодействия
gain--Weak-error и два эффекта коррекции перестали быть значимыми.
Все главные примеры статьи значимы в обоих расчётах.
Первоначальная и уточнённая таблицы сохранены полностью, включая
неудобные для интерпретации граничные случаи.

\subsection{Сводка всей области проверки}
\begin{table}[htbp]\centering\small
\caption{Полный контекст выбранных иллюстраций. Числа --- количества
контрастов после общей поправки своего этапа. Для C указаны
первичные результаты; отличия уточнения описаны выше.}
\label{tab:scope_audit}
\begin{tabular}{lp{.48\linewidth}r}
\toprule
Семья & Вид сравнения & Значимых / всего \\
\midrule
B & Простые accuracy-различия: natural & 0 / 756 \\
B & Простые accuracy-различия: rescue8 & 0 / 756 \\
B & Простые accuracy-различия: harm8 & 77 / 756 \\
B & Взаимодействия метода и состава & 56 / 1\,512 \\
B & В том числе gain-only взаимодействия & 0 / 882 \\
B & Коррекция IPW--raw & 34 / 336 \\
B & Natural-профили спасений/вреда & 0 / 1\,512 \\
B & Selector--confidence & 0 / 63 \\
\midrule
C & Gain-only взаимодействия, accuracy-настройка & 8 / 126 \\
C & Gain-only взаимодействия, MSE-настройка & 0 / 126 \\
C & Прямые эффекты смены tuning-критерия & 0 / 144 \\
C & Влияние критерия на gain--Weak-error interaction & 0 / 42 \\
\midrule
M & Accuracy, natural & 1 / 1\,188 \\
M & Macro-precision, natural & 0 / 1\,188 \\
M & Macro-F1, natural & 1 / 1\,188 \\
\bottomrule
\end{tabular}
\end{table}

Полные CSV содержат все 5\,691, 1\,698 и 3\,564 контраста,
а не только строки этой сводки. На BANKING77 28 raw-ухудшений
относительно natural уже входят в семью B как взаимодействия
с confidence; нового семейства для них не добавлялось.
Четыре простых gain-only различия harm8 в B относятся к одной
low-label паре TF10 $\to$ TFfull. Прежние двухметодные контроли
с 30 seeds и отдельной семьёй 14 сравнений сохранены в исходной
версии статьи; их уровни значимости не подменяют общую проверку B.

\section{Контрпример: агрегированной таблицы ошибок недостаточно}
Пусть два значения признаков $z$ равновероятны, а столбцы
таблицы задают условные вероятности совместных исходов.
\begin{table}[htbp]\centering\small
\caption{Популяции с одинаковыми агрегатами, но разной реакцией
идеального gain-ранжирования на harm-отбор.}
\begin{tabular}{ccrrrr}
\toprule
Популяция & $z$ & $P(0,0\mid z)$ & $P(0,1\mid z)$ & $P(1,0\mid z)$ & $P(1,1\mid z)$ \\
\midrule
A & 1 & 0.02 & 0.08 & 0.02 & 0.88 \\
A & 2 & 0.06 & 0.24 & 0.06 & 0.64 \\
B & 1 & 0.02 & 0.08 & 0.08 & 0.82 \\
B & 2 & 0.06 & 0.24 & 0.00 & 0.70 \\
\bottomrule
\end{tabular}
\end{table}

У обеих популяций агрегированная таблица одинакова:
$(0.04,0.16,0.04,0.76)$, как и ошибка первого эксперта
при каждом $z$: 0.1 и 0.3. Исходный gain выбирает $z=2$.
При $a=8$ популяция A имеет
$g_{a,1}=(0.08-8\cdot0.02)/(1+7\cdot0.02)\approx-0.0702$
и $g_{a,2}=(0.24-8\cdot0.06)/(1+7\cdot0.06)\approx-0.1690$,
поэтому порядок обращается. В B получаются
$g_{a,1}\approx-0.3590$ и $g_{a,2}=0.24$, порядок сохраняется.
Итак, совпадение accuracy экспертов и глобального rescue/harm ratio
не определяет реакцию ранжирования. Пример характеризует
идеальные условные цели, а не новые обученные модели на реальном тесте.

\section{Происхождение данных и воспроизводимость}
\paragraph{Источники и параметры.}
BANKING77: PolyAI revision
\texttt{57ec275d8078af65b7731c2a98be812d844a6d6b}, CC BY 4.0.
CLINC150: \texttt{clinc/oos-eval}, revision
\texttt{828f8093932c8fe6ca7936c3d2e52903b1c523de},
\texttt{data/data\_full.json}. Нормализация CLINC использует
NFKC, casefold и свёртку пробелов; конфликты train-меток удаляются,
одинаковые train-тексты одной метки объединяются. Dev очищается
от повторов, конфликтов и пересечений с train; official test сохраняется.

TF--IDF использует word n-grams (1,2), min\_df=2,
sublinear TF и не более 30\,000 признаков. Полный LR выбирает
$C\in\{0.25,1,4,16\}$; frozen MiniLM --- нормализованный mean pooling
и LR с $C\in\{0.25,1,4,16,64\}$.
Pretrained \texttt{sentence-transformers/all-MiniLM-L6-v2}
имеет revision \texttt{1110a243fdf4706b3f48f1d95db1a4f5529b4d41}.
BANKING77 BERT-tiny использует snapshot
\texttt{6f75de8b60a9f8a2fdf7b69cbd86d9e64bcb3837},
mean pooling, линейную голову, восемь эпох refit с AdamW,
batch32, encoder LR $10^{-4}$, head LR $10^{-3}$,
weight decay0.01, warmup/linear decay, gradient norm cap1,
длину до64токенов и seed20261002.
Прогнозы всех семи BANKING77 экспертов, включая tuned MiniLM,
берутся из зафиксированных архивов; полные параметры и хеши
приведены в соответствующих manifests предыдущего этапа.

\paragraph{Аудит вычислений.}
Полный BANKING77 replay воспроизводит 110\,880 выбранных моделей
и 763\,560 строк метрик; CLINC replay --- 23\,040 кандидатов
и 79\,200 строк. Многоклассовый этап использует эти же решения
и исходные вероятности экспертов, выравнивает row IDs и class labels,
проверяет соответствие бинарной правильности $W,S$ и включает
переходы между двумя разными неправильными классами.
При fit1000/20\%/official сохранены TP и predicted-support
каждого класса и каждого seed; средние классовых метрик вычисляются
по запускам, а не по усреднённой матрице ошибок.



\paragraph{Материалы и воспроизводимость.}
Исходная статья и базовый эксперимент сохраняются без изменения текста.
Текущие исходники, результаты сборки и презентация находятся в папке
\texttt{main/}; файлы \texttt{main.tex} и \texttt{main.pdf} относятся к статье,
\texttt{slides.pdf} --- к презентации.
Новые экспериментальные разделы находятся в \texttt{additions/}.
Полные контрасты, настройки и метрики сохранены в отдельных
исследовательских архивах BANKING77, CLINC150 и multiclass-анализа.
Числа основных таблиц перенесены из проверенных CSV с сохранением
полных статистических семей и направлений сравнений.
Файл \texttt{verification.json} фиксирует контроль исходной аннотации,
прежних разделов и входных артефактов.
